\documentclass[11pt]{article}
\usepackage[a4paper,margin=25mm]{geometry}
\usepackage{amsmath,amssymb,amsthm}
\usepackage[T1]{fontenc}
\usepackage{lmodern}
\usepackage[hidelinks]{hyperref}
\newtheorem{theorem}{Theorem}
\newcommand{\Q}{\mathbb Q}
\title{Jacobson Radical Powers Need Not Vanish\newline at Any Finite Stage in an FBN Ring}
\author{Conjecture 00000009978}
\date{}
\begin{document}
\begin{center}
{\Large Jacobson Radical Powers Need Not Vanish\\
at Any Finite Stage in an FBN Ring}\par\medskip
Conjecture 00000009978
\end{center}
\begin{abstract}
The finite-termination clause of Conjecture 00000009978 is false. The formal power series ring $R=\Q[[X]]$ is fully bounded Noetherian, its Jacobson radical is $J=(X)$, and $J^N\ne0$ for every $N\geq0$. Its finite power intersections therefore never vanish, although its infinite power intersection is zero. The distinction is proved both mathematically and in Lean using actual formal power series and ideals.
\end{abstract}

\section*{The clause and its precise meaning}
The original English statement asserts that the intersection of Jacobson radical powers in FBN rings ``vanishes in finitely many steps''; the Chinese statement likewise says that the intersection is zero at a finite step. Neither version assumes that the ring is Artinian or finite dimensional over a field.

Write $J=J(R)$ for the Jacobson radical, the intersection of maximal ideals in the commutative case. The finite stages of the descending power intersection are
\[
 F_N=\bigcap_{j=0}^{N}J^j=J^N,\qquad N\geq0,
\]
because ideal powers descend and $J^0=R$. Thus the finite-termination clause requires $J^N=0$ for some finite $N$. Omitting $J^0$ gives the same intersection for $N\geq1$.

\begin{theorem}
For the fully bounded Noetherian ring $R=\Q[[X]]$,
\[
 J(R)=(X),\qquad J(R)^N=(X^N)\ne0\quad (N\geq0).
\]
The powers descend strictly at every stage, while
\[
 \bigcap_{N\geq0}J(R)^N=0.
\]
In particular, $F_N\ne0$ for every finite $N$, disproving the finite-termination clause.
\end{theorem}
\begin{proof}
Every nonzero formal power series $f$ has a least nonzero coefficient, say in degree $r$, and factors as $f=X^r u$ with $u$ a unit. A series is a unit exactly when its constant coefficient is nonzero. For a nonzero ideal $I$, choose the least order $r$ among its nonzero elements. The factorization shows $X^r\in I$, and minimality shows every element of $I$ is divisible by $X^r$. Hence $I=(X^r)$. Every ideal is therefore principal, so $R$ is Noetherian.

The ring is also FBN: every essential right ideal of each prime quotient must contain a nonzero two-sided ideal \cite[p.~1]{cg}. Here ``essential'' means intersecting every nonzero right ideal nontrivially. Each prime quotient is a nonzero commutative ring, so its left, right, and two-sided ideals coincide. An essential ideal is nonzero by intersecting it with the whole ring, and contains the required two-sided ideal, namely itself.

The nonunits of $R$ are exactly the series with zero constant coefficient, which form $(X)$. Therefore $(X)$ is the unique maximal ideal and $J(R)=(X)$. It follows that $J(R)^N=(X^N)$. The coefficient of degree $N$ in $X^N$ is $1$, so $X^N\ne0$ and $J(R)^N\ne0$. Moreover, every element of $(X^{N+1})$ has coefficient zero in degree $N$. Thus $X^N\notin J(R)^{N+1}$, proving strict descent.

Finally, if $f$ lies in every $J(R)^N$, then for each $r\geq0$ it lies in $(X^{r+1})$. Its coefficient in degree $r$ is consequently zero. All its coefficients vanish, so $f=0$ and the infinite intersection is zero.
\end{proof}

\section*{Scope of the disproof}
This refutes the source's finite-termination conjunct for a genuine FBN ring. It does not refute the zero \emph{infinite} intersection, which is proved for the same ring, or concern the nilpotence of radicals of Artinian rings. No numerical value or interpretation is assigned to the source's unspecified exponential decay constant.

\section*{Lean formalization}
The file \texttt{lean/Main.lean} uses Lean 4.19.0 and Mathlib at commit
\begin{center}\small\texttt{c44e0c8ee63ca166450922a373c7409c5d26b00b}.\end{center}
The correspondence with the mathematical proof is explicit.
\begin{itemize}
\item \texttt{Ring} is Mathlib's actual \texttt{PowerSeries $\Q$}, with untruncated coefficients and formal-series ring operations.
\item \texttt{comm\_noetherian\_fbn} proves \texttt{CommFBN} for every commutative Noetherian ring by checking essential ideals of prime quotients.
\item \texttt{radical} is Mathlib\textquotesingle s \texttt{Ideal.jacobson} of the zero ideal; local-ring and power-series theorems identify it with $(X)$.
\item \texttt{radical\_pow\_nonzero} and \texttt{radical\_powers\_strictly\_descend} hold for every natural $N$, witnessed by the coefficient of $X^N$.
\item \texttt{prefixIntersection} is the actual finite ideal intersection indexed by \texttt{Fin (N+1)}. Its equality with $J^N$ and nontermination are proved.
\item \texttt{infinite\_intersection\_zero} proves the zero infinite intersection by coefficient extensionality. \texttt{counterexample} combines this with FBN and finite-stage nontermination.
\end{itemize}
The FBN encoding specializes to commutative rings, where right ideals are precisely the ordinary ideals used in the file; their right-multiplication closure is also proved. This exactly covers the exhibited ring. Every stage is handled by a universally quantified theorem, with no numerical approximation or bounded search.

\section*{Source and verification}
\begin{flushleft}
\texttt{SOURCE.md} preserves the original
\href{https://github.com/The-Last-Math-Competition/The-Last-Math-Competition/blob/main/conjectures/00000009978.md}{conjecture 00000009978} verbatim, and \texttt{README.md} gives build instructions.
The \texttt{verification} directory contains the provenance record
\texttt{SOURCE\_PROVENANCE.md}, the report \texttt{BUILD.json}, and the build, axiom-audit, and PDF logs.
\end{flushleft}

\begin{thebibliography}{1}
\bibitem{cg} S.~Caenepeel and T.~Gu\'ed\'enon,
\emph{Fully bounded noetherian rings and Frobenius extensions},
arXiv:math/0503686v3, 2006.
\url{https://arxiv.org/abs/math/0503686v3}.
\end{thebibliography}
\end{document}
